\documentclass[12pt]{article}
\usepackage[utf8]{inputenc}
\usepackage[spanish]{babel}
\usepackage{graphicx}
\usepackage{float}
\usepackage{caption}
\usepackage{csquotes}
\usepackage{booktabs}
\usepackage{amsmath}
\usepackage{hyperref}
\usepackage[backend=biber,style=apa]{biblatex}
\addbibresource{../referencias.bib}
\usepackage{geometry}
\usepackage{url}
\geometry{margin=2.5cm}

\begin{document}
\begin{titlepage}
\centering

\vspace*{\fill}
{\Large \textbf{Proyecto CA0305}} \\
{\large Bitácora 1} \\
\vspace{2cm}
{\large \textbf{Integrantes}} \\
\vspace{1cm}
\begin{tabular}{p{6cm} p{2cm}}
Ashly Garro Villanueva & C33198 \\
Robert Meléndez Jiménez & C34772 \\
Adriana Monge Rodríguez & C35042 \\
Emily Sánchez Mancía & C27260 \\
Alessandro Umaña Vega & C37963 \\
\end{tabular}

\vspace{2cm}

{\large \textbf{I Ciclo 2026}}
\vspace*{\fill}

\end{titlepage}
\tableofcontents
\newpage

% -------------------------
\section{Fundamentación Matemática}
\subsection{Linear Layer Forward}

\subsection{Linear Layer Backward}

\subsection{ReLU Activation Forward}

\subsection{ReLU Activation Backward}

\subsection{Sigmoid Activation}

\subsection{Tanh Activation}

\subsection{MSE Loss Forward \& Backward}

Lo que hacemos en este ejercicio es crear una función que recibe como parámetros los valores de nuestra predicción y los valores reales que debieron obtenerse. Primero, se comparan ambos valores y se calcula la diferencia entre ellos, la cual representa el error. Para evitar que el error sea negativo y para penalizar más los errores grandes, este error se eleva al cuadrado.

Como queremos obtener el promedio del error cuadrático, se utiliza `.mean()` para calcular la pérdida, también llamada `loss`.

MSE, por sus siglas en inglés *Mean Squared Error*, es una de las funciones de pérdida más utilizadas en machine learning. Esta función mide la diferencia entre las predicciones y los valores reales, y es especialmente útil porque castiga con mayor fuerza los errores grandes en comparación con los errores pequeños.

También es importante calcular el gradiente, ya que este indica cómo cambia la pérdida con respecto a cada valor predicho. La fórmula de la pérdida MSE es:

$$
L = \frac{1}{N \cdot D} \sum_{k,l} \left( y_{\text{pred},kl} - y_{\text{true},kl} \right)^2
$$

Para derivar con respecto a un valor específico $(y_{\text{pred},ij})$, 
se plantea:

$$
\frac{\partial L}{\partial y_{\text{pred},ij}} =
\frac{\partial}{\partial y_{\text{pred},ij}}
\left[
\frac{1}{N \cdot D}
\sum_{k,l}
\left(
y_{\text{pred},kl} - y_{\text{true},kl}
\right)^2
\right]
$$

Como solo el término ((i,j)) aparece en esa derivada, tq:

$$
\frac{\partial L}{\partial y_{\text{pred},ij}} =
\frac{1}{N \cdot D}
\cdot
2
\left(
y_{\text{pred},ij} - y_{\text{true},ij}
\right)
$$

Por lo tanto, el gradiente:

$$
\frac{\partial L}{\partial y_{\text{pred},ij}} =
\frac{2}{N \cdot D}
\left(
y_{\text{pred},ij} - y_{\text{true},ij}
\right)
$$

El número 2 aparece porque al derivar un término elevado al cuadrado, baja el exponente. Además, el factor $(\frac{1}{N \cdot D})$ normaliza el resultado según la cantidad total de elementos.

En Python, esto se representa calculando primero el error:

$$error = y_{pred} - y_{true}$$

Luego, la pérdida se obtiene con el promedio del error al cuadrado:

loss = np.mean(error ** 2)

Y el gradiente se calcula aplicando la fórmula:

dx = (2 / error.size) * error

La idea clave es que el gradiente es proporcional al error: cuando el error es mayor, el gradiente también será mayor. Esto permite saber cuánto debe ajustarse la predicción para reducir la pérdida. Sin embargo, en los métodos de optimización, se avanza en la dirección contraria al gradiente, porque el gradiente señala hacia dónde aumenta la pérdida.

\subsection{Binary Cross Entropy Loss}

\subsection{Categorical Cross Entropy Backward}
En este ejercicio se crea una función llamada categorical\_ce\_backward, la cual recibe como parámetros probs y targets. La variable probs contiene las probabilidades generadas por la función softmax para cada muestra y cada clase, mientras que targets contiene la clase correcta de cada muestra. 
Entonces, ambos valores se convierten en arreglos de NumPy para poder trabajar con ellos de forma vectorizada. Luego, se obtiene N usando probs.shape[0], ya que N representa la cantidad de filas, es decir, la cantidad de muestras. 
Después, se crea una copia de probs llamada dx, para no modificar los datos originales. Como el gradiente de softmax combinado con entropía cruzada se simplifica a probs y\_true lo que se hace es restar 1 en la posición correspondiente a la clase correcta de cada fila. Esto se logra con $dx[np.arange(N), targets] -= 1$. 
Como ultimo, se divide dx entre N porque el ejercicio pide obtener el gradiente promedio. La función devuelve dx, que representa el gradiente con respecto a los logits. 

La pérdida utilizada es la combinación de Softmax + Cross Entropy. Primero, la función softmax convierte los logits $(z_{ij})$ en probabilidades $(p_{ij})$, donde cada valor representa la probabilidad de que la observación (i) pertenezca a la clase (j).

La pérdida Cross Entropy toma únicamente la probabilidad de la clase correcta $(p_{i,y_i})$ y calcula su logaritmo negativo. Por eso, la pérdida se escribe como:

$$
L = -\frac{1}{N}\sum_{i=1}^{N}\log(p_{i,y_i})
$$

Al derivar esta pérdida con respecto al logit $(z_{ij})$, se deben considerar las derivadas de softmax y de Cross Entropy. Aunque al inicio aparecen varios términos, incluyendo términos cruzados entre clases, estos se simplifican.

Después de aplicar la regla de la cadena y simplificar, el resultado final queda como:

$$
\frac{\partial L}{\partial z_{ij}} = p_{ij} - y_{ij}
$$

donde $(y_{ij})$ vale 1 si la clase (j) es la clase correcta, y 0 si no lo es.

En Python, esto se representa copiando la matriz de probabilidades y restando 1 únicamente en la posición de la clase correcta de cada fila:

$$
dx = probs
$$

$$
dx[i, y_i] = dx[i, y_i] - 1
$$

Finalmente, como la pérdida se calcula como un promedio para las (N) observaciones, el gradiente se divide entre (N):

$$
dx = \frac{dx}{N}
$$

La idea clave es que la derivación completa de Softmax + Cross Entropy se reduce a una forma muy simple: a 
las probabilidades predichas se les resta la matriz real codificada en one-hot.


\subsection{SGD Optimizer}
\subsection{Momentum Optimizer}

\subsection{RMSProp Optimizer}

Este ejercicio trata sobre implementar la función rmsprop\_step(w, dw, lr, decay, eps, cache), la cual se utiliza para actualizar los pesos de un modelo mediante el optimizador RMSProp.

RMSProp es un método de optimización que ajusta la tasa de aprendizaje de forma adaptativa. Esto significa que no actualiza todos los pesos de la misma manera, sino que toma en cuenta el comportamiento reciente del gradiente para hacer actualizaciones más estables.

El primer paso consiste en actualizar el caché:

$$
s_{\text{nuevo}} = \beta \cdot s + (1 - \beta) \cdot dw^2
$$

En esta fórmula, (s) representa el caché anterior, (dw) es el gradiente actual y $(\beta)$, también llamado decay, controla cuánta memoria se conserva de los gradientes anteriores. Este caché funciona como un promedio móvil de los gradientes al cuadrado.

Luego, se calcula una tasa de aprendizaje adaptativa:

$$
\text{aprendizaje adaptativo} =
\frac{lr}{\sqrt{s_{\text{nuevo}}} + \epsilon}
$$

Aquí, (lr) representa la tasa de aprendizaje original y $(\epsilon)$ es un valor muy pequeño que se utiliza para evitar divisiones entre cero. La raíz cuadrada del caché permite ajustar el tamaño del paso dependiendo de qué tan grandes han sido los gradientes.

Después, se actualizan los pesos:

$$
w_{\text{nuevo}} = w - \text{aprendizaje adaptativo} \cdot dw
$$

Esto significa que los pesos se mueven en la dirección contraria al gradiente, ya que el objetivo es reducir la pérdida del modelo.

La función finalmente retorna dos valores: los pesos actualizados $(w_{\text{nuevo}})$ y el nuevo caché $(s_{\text{nuevo}})$. Este nuevo caché se debe guardar porque será utilizado en la siguiente iteración, por eso RMSProp es un proceso iterativo.

La tasa de decaimiento b controla la memoria del algoritmo. Si b = 0, solo se toma en cuenta el gradiente actual. Si b = 0.9, se conserva el 90\% del caché anterior, que suele ser un valor común. Si b = 0.99, el algoritmo conserva más memoria, por lo que la adaptación es más lenta pero más estable.

En resumen, RMSProp ayuda a que el entrenamiento sea más estable porque divide el gradiente entre una medida acumulada de los gradientes recientes. De esta forma, evita actualizaciones demasiado grandes cuando los gradientes son altos y permite un mejor control del aprendizaje.

\subsection{Adam Optimizer}

\subsection{Weight Initialization (Xavier/Glorot)}
La inicialización Xavier fue propuesta por \textcite{GLOROT2010} para 
mantener aproximadamente constante la varianza de las activaciones y de 
los gradientes a través de las capas de una red neuronal profunda. 
Esto ayuda a evitar que las señales se amplifiquen o atenúen excesivamente 
durante la propagación hacia adelante y hacia atrás, facilitando el 
entrenamiento de la red.

Los autores observaron que, cuando la varianza de las señales cambia 
significativamente entre capas, el entrenamiento se vuelve difícil 
debido a que las activaciones y los gradientes pueden crecer o 
disminuir progresivamente, donde las activaciones corresponden a las salidas 
producidas por las neuronas después de aplicar la función de activación, 
mientras que los gradientes representan la sensibilidad de la función de 
pérdida respecto a los parámetros de la red.

Considere una capa lineal definida por $y = Wx+b$ donde $x$ representa las 
entradas, $W$ la matriz de pesos y $b$ el vector 
de sesgos. 

Si asumimos que las entradas son independientes y tienen varianza $var(x)$, 
la varianza de la salida $y$ puede aproximarse por:
$$var(y) \approx fan_{in} \cdot var(W)var(x)$$ donde $fan_{in}$ es la 
cantidad de conexiones de cada neurona. Esto nos muestra que la 
varianza de la salida depende directamente del número de entradas 
y de la varianza de los pesos.

Para que la señal conserve aproximadamente la misma magnitud al pasar de 
una capa a otra, se busca que $var(y) \approx var(x)$, esto implica
que $fan_{in} \cdot var(W) = 1$. De esta manera, la varianza 
de los pesos debe ser inversamente proporcional al número de entradas.

Sin embargo, \textcite{GLOROT2010} señalan que también es necesario 
considerar la propagación hacia atrás de los gradientes durante el 
proceso de retropropagación. Si los gradientes aumentan o disminuyen 
excesivamente al atravesar las capas, las actualizaciones de los pesos 
se vuelven inestables o ineficientes. Para equilibrar tanto la propagación 
hacia adelante como la retropropagación, los autores incorporan además 
el número de conexiones de salida, denotado por $fan_{out}$.

A partir de este análisis, proponen inicializar los pesos
mediante una distribución uniforme simétrica alrededor de cero
cuya varianza sea 
$$Var(W)=\frac{2}{fan_{in}+fan_{out}}.$$
Dado que la varianza de una distribución uniforme ($U(-a,a)$) es $Var(W)=\frac{a^2}{3}$
al igualar ambas expresiones se obtiene 
$$a=\sqrt{\frac{6}{fan_{in}+fan_{out}}}.$$ 

Por lo tanto, los pesos se inicializan según
$$W \sim U(-\sqrt{\frac{6}{fa_{in}+fan_{out}}}, \sqrt{\frac{6}{fa_{in}+fan_{out}}})$$

La idea central de este método es mantener una magnitud relativamente 
estable de las activaciones y de los gradientes a través de las capas 
de la red. De esta forma, se reduce el riesgo de desvanecimiento o 
explosión de gradientes, permitiendo un entrenamiento más estable y eficiente.

\subsection{Weight Initialization (Kaiming/He)}

\subsection{Dropout Forward}

\subsection{Dopout Backward}

El fin de este ejercicio es crear la funcion dropout\_backward(dout, mask, p, train = True), que es la que calcula el gradiente, hay 2 casos si se encuentra en entrenamiento o no. Matemáticamente, durante el entrenamiento se utiliza una máscara binaria (m), donde cada valor puede ser 0 o 1. Si m = 0, la neurona se apaga y no participa en el cálculo; si m = 1, la neurona permanece activa.

En el forward pass, usando dropout invertido, la salida se calcula como:

$$
out = \frac{x \cdot m}{1-p}
$$

donde (p) representa la probabilidad de apagar una neurona, y (1-p) representa la probabilidad de mantenerla activa. La división entre (1-p) se hace para compensar las neuronas apagadas y mantener estable la escala de los valores durante el entrenamiento.

Luego, en el backward pass, se aplica la regla de la cadena. Como:

$$
out = \frac{x \cdot m}{1-p}
$$

entonces la derivada de (out) con respecto a (x) es:

$$
\frac{\partial out}{\partial x} = \frac{m}{1-p}
$$

Por eso, el gradiente queda como:

$$
dx = dout \cdot \frac{m}{1-p}
$$

dx = mask * dout / (1-p)

Si una posición de la máscara vale 0, el gradiente en esa neurona también será 0, porque esa neurona fue apagada. Si la máscara vale 1, el gradiente pasa, pero se escala dividiéndolo entre (1-p).

Cuando el modelo no está en entrenamiento, no se aplica dropout, porque en evaluación se utilizan todas las neuronas. Por eso, en ese caso el gradiente simplemente se devuelve igual:

dx = dout

En resumen, la fundamentación matemática del ejercicio se basa en la regla de la cadena. El gradiente que viene de la capa siguiente, (dout), se multiplica por la derivada de la operación dropout. Por eso, durante entrenamiento el gradiente se filtra con la máscara y se escala por (1-p), mientras que durante evaluación se mantiene sin cambios.

\subsection{Batch Normalization Forward}

\subsection{Batch Normalization Backward}

\subsection{Layer Normalization Forward \& Backward}
\newpage
\printbibliography
\end{document}